%==============================================================================
% Appendix for Objective 4
% Everything below this line is appendix material. Chapters are lettered
% (A, B, ...) and sections numbered A.1, A.2, ... Reference from the body text
% with e.g. Appendix~\ref{app:obj4} or Appendix~\ref{app:obj4-code}.
%==============================================================================
\clearpage
\appendix
\renewcommand{\chaptermark}[1]{\markboth{Appendix \thechapter: #1}{}}

\chapter{Objective 4: Supporting Material}
\label{app:obj4}

\section{SSH Setup}
\label{app:obj4-ssh-setup}

\begin{lstlisting}[language=bash,caption={The commands used to setup SSH connection between the Raspberry Pi and a remote computer.},label={lst:ssh-setup}]
sudo apt install openssh-server
sudo systemctl enable ssh
sudo ufw allow ssh
\end{lstlisting}

\section{GPIO Packages Install}
\label{app:obj4-gpio-packages}

\begin{lstlisting}[language=Bash,caption={Bash command used to install python \gls{gpio} library.},label={lst:gpio-package-install}]
sudo apt install python3-lgpio gpiod
\end{lstlisting}

\section{Virtual Environment and \texttt{rpi-lgpio} Install}
\label{app:obj4-rpi-lgpio-install}

\begin{lstlisting}[language=Bash,caption={Commands used to create the python virtual environment and install \texttt{rpi-lgpio}.},label={lst:rpi-lgpio-install}]
wget http://abyz.me.uk/lg/lg.zip
unzip lg.zip
cd lg
make
sudo make install
python3 -m venv gpio-env
source gpio-env/bin/activate
pip install rpi-lgpio
\end{lstlisting}

\section{GPIO Read Script}
\label{app:obj4-gpio-script}

\begin{lstlisting}[language=python,caption={Python script used to run on the Raspberry Pi to monitor one of the \gls{gpio} pins.},label={lst:gpio-script}]
#!/usr/bin/env python3
import RPi.GPIO as GPIO # Uses rpi-lgpio backend automatically
from time import sleep
import signal

# Setup
GPIO_PIN = 17
GPIO.setmode(GPIO.BCM)
GPIO.setup(GPIO_PIN, GPIO.IN, pull_up_down=GPIO.PUD_DOWN)

print("Sensing 3.3V High...")

try:
    while True:
        if GPIO.input(GPIO_PIN):
            print("Pin High")
            sleep(1)
except KeyboardInterrupt:
    print("Stopping...")
finally:
    GPIO.cleanup()
pause()
\end{lstlisting}



\section{Telemetry Data Visualisation Script}
\label{app:obj4-telem-data-visualisation}

\begin{lstlisting}[language=python,caption={Python script used to plot and visualise syncrhonised radar and telemetry data},label={lst:telem-vis-script}]
import folium
import pandas as pd
import os

# parameters
filename = "flight_log_5.csv"


# Code
df = pd.read_csv(filename)

# first lat/lon sample is invalid, drop it and reindex from 0
df = df.iloc[1:].reset_index(drop=True)

output_file = os.path.splitext(filename)[0] + ".html"  # -> "flight_log_4.html"
os.remove("synced_data.html") if os.path.exists("synced_data.html") else None

m = folium.Map(location=[df.at[0, "latitude"], df.at[0, "longitude"]], zoom_start=100)

# Create list of [latitude, longitude] points
points = list(zip(df["latitude"], df["longitude"]))

# Draw flight path
folium.PolyLine(points, weight=3).add_to(m)

capture_count = 0
for i, row in df.iterrows():
    tooltip_html = (
        f"<b>#{i}</b><br>"
        f"Latitude: {row['latitude']:.6f}<br>"
        f"Longitude: {row['longitude']:.6f}<br>"
        f"Capture Count: {row['radar_capture_count']}<br>"
    )

    if row["radar_capture_count"] == capture_count + 1:
        folium.CircleMarker(
            location=[row["latitude"], row["longitude"]],
            radius=3,
            color="blue",
            fill=True,
            fill_opacity=0.8,
            tooltip=tooltip_html,
        ).add_to(m)
        capture_count = 1 + capture_count
# Save map
m.save(output_file)

\end{lstlisting}